%% ============================================================
%% Chapter 14: Model Capsules and Ternary Learning
%% ============================================================

\chapter{Model Capsules and Ternary Learning}
\label{ch:ch14-model-capsules}
\section{The Deployment-Gap Problem}

Chapter~\ref{ch:ch13-verification-inheritance} treated the question of whether a claim survives its own transformation, formally and in general. This chapter treats a specific, concrete instance of the same question, closer to the machine than to the document: whether a trained model's deployed behavior survives the transformations, quantization, export, kernel substitution, tokenizer binding, that separate optimization from execution.

A conventional model pipeline often contains several non-identical machines. Optimization uses floating-point shadow weights and training kernels; export applies quantization and layout conversion; inference invokes a separate runtime; and deployment adds platform-specific tokenization, scaling, memory, and scheduling behavior. Every transformation can preserve average benchmark performance while changing individual outputs, numerical stability, latency, or failure behavior. The relevant object is therefore not merely a parameter tensor $\theta$, but the complete deployed computation,
\[
F = (G, Q, S, T, K, R, P),
\]
where $G$ is the graph or bytecode, $Q$ the weight alphabet and packing rule, $S$ the scale metadata, $T$ the tokenizer and preprocessing state, $K$ the kernel semantics, $R$ the runtime contract, and $P$ the execution policy. Two artifacts with identical ternary weights can still implement different functions when any other component differs.

Deployment-native training does not require every training operation to be low precision. Gradient accumulation and optimizer state may remain floating point. It requires that the forward operator used to estimate the task loss faithfully reproduce the discrete weights, scales, rounding, saturation, activation treatment, and layer ordering that will be used after export.

Deployment drift separates into components arising at distinct stages of the pipeline, each admitting a different remedy. Representation drift arises when the forward operator used for loss estimation differs from the exported operator, most commonly through a mismatched straight-through surrogate or an export-time rounding rule not exercised in training. Kernel drift arises when the mathematically specified operator is realized by different low-level implementations across devices, for example a fused kernel that reorders accumulation and changes rounding under finite precision. Preprocessing drift arises when the tokenizer, normalization, or templating logic bound to the training data differs from the version bound to the runtime capsule, which can change token boundaries without changing any weight. Scheduling drift arises when batching, padding, or caching strategies at inference time alter numerical results relative to the unbatched computation used during evaluation. These four sources are not mutually exclusive and can partially cancel or compound; a pipeline that reports only aggregate deployment disagreement without attributing it to a source cannot say which remedy would help.

\section{Ternary Parameterization}

Let a layer maintain real-valued shadow parameters $W \in \mathbb{R}^{m\times n}$ while executing a ternary forward matrix $\hat{W}$. A simple symmetric quantizer uses a scale $\alpha > 0$ and threshold $\Delta \geq 0$:
\[
\hat{W}_{ij} = \alpha \cdot q(W_{ij}/\alpha), \qquad q(x) \in \{-1,0,+1\}, \qquad q(x) = \begin{cases} -1 & x < -\Delta \\ 0 & |x| \leq \Delta \\ +1 & x > \Delta. \end{cases}
\]
The ternary symbol requires $\log_2 3 \approx 1.585$ bits of information under a uniform code, although practical packing may use two physical bits per weight or block encodings with amortized overhead; calling the representation strictly one-bit is therefore imprecise. Its computational appeal is that multiplication by a ternary weight reduces to signed accumulation and omission, subject to the cost of scales, activation arithmetic, packing, memory movement, and kernel dispatch. Training commonly uses a straight-through approximation because $q$ is piecewise constant: the backward pass substitutes a surrogate derivative for $\partial q/\partial W$, which means the learned solution depends on the surrogate, clipping rule, scale estimator, and update schedule, so reproducibility requires recording these rules, not only the final packed weights.

The symmetric, per-tensor quantizer is the simplest member of a larger family, and the family member chosen is itself part of the deployment specification $F$ rather than an incidental training detail. A per-channel quantizer assigns an independent scale $\alpha_k$ to each output channel, reducing quantization error on layers with heterogeneous weight magnitude at the cost of storing a vector of scales rather than a scalar; the capsule manifest must then record the reduction axis and count, since a runtime that assumes a single scalar scale will silently misinterpret a per-channel capsule rather than failing loudly. An asymmetric variant introduces an offset $\beta$, producing reconstruction levels $\{\beta-\alpha,\beta,\beta+\alpha\}$, which is a three-level operator rather than a zero-centered ternary weight in the sense used elsewhere: signed accumulation eliminates multiplication only when the reconstruction levels are exactly $\{-1,0,+1\}$, so a deployment adopting the asymmetric variant should report it as a distinct operator with its own efficiency profile.

Deployment equivalence itself needs a measurable definition. Let $f_{\mathrm{train}}(x;\theta)$ denote the forward pass used while training and $f_{\mathrm{run}}(x;C)$ the runtime evaluation of a serialized capsule $C$. For a test distribution $D$,
\[
\delta_{\mathrm{out}} = \Pr_{x\sim D}\bigl[ \arg\max f_{\mathrm{train}}(x;\theta) \neq \arg\max f_{\mathrm{run}}(x;C) \bigr],
\]
\[
\delta_{\mathrm{num}} = \mathbb{E}_{x\sim D}\left[ \frac{\|f_{\mathrm{train}}(x;\theta) - f_{\mathrm{run}}(x;C)\|_2}{\|f_{\mathrm{train}}(x;\theta)\|_2 + \varepsilon} \right].
\]
Exact bitwise agreement may be unrealistic across heterogeneous devices, but token-level disagreement, normalized logit error, and task-level divergence are measurable, and a pipeline should state which equivalence criterion it promises and on which platforms it has been tested.

\section{The Model Capsule}

The deployable unit should be a self-describing capsule rather than an unstructured weight file, binding executable semantics to model state and supporting hashing, signature verification, compatibility checks, and deterministic reconstruction. Its required components are a manifest recording format version, model identity, hashes, license, and provenance; a program, typed bytecode or a restricted graph specifying computation independently of host-language objects; parameters, packed ternary weights and scale blocks preserving the executed representation; a text interface, tokenizer, normalization, vocabulary, and templates, avoiding silent input mismatch; a runtime contract specifying operators, numerical rules, memory limits, and ABI, making compatibility testable; and evidence, an evaluation summary, calibration set hash, and build attestation recording what has actually been measured. A capsule should not contain arbitrary native code. A restricted instruction set with bounded tensor shapes, declared memory, immutable parameter regions, and explicit operator versions is easier to validate and isolate. The signature authenticates the bytes and declared producer; it does not establish model quality or safety.

A capsule format that changes without a disciplined versioning rule reintroduces the deployment gap it was built to close, since a runtime and a capsule produced under different silent assumptions can both parse successfully while disagreeing on semantics. A workable convention assigns three independent version fields to the manifest: a format version governing the container schema itself, a contract version governing the runtime operators, numerical rules, and ABI, and a content version identifying the specific trained artifact. A conservative runtime may require exact format-version equality. Contract versions should be compared under a declared compatibility relation, for example a runtime advertising contract version $v$ accepting any capsule declaring a contract version in a closed range $[v_{\min}, v]$ that the runtime publishes, so older capsules remain executable under runtimes that have only added operators and newer capsules are refused rather than silently misexecuted. Content versions carry no compatibility semantics of their own; they exist so that two capsules with identical format and contract versions can still be distinguished, hashed, and compared for provenance. Refusal is warranted whenever the contract-version range test fails; acceptance is not warranted merely because the capsule parses. This three-field versioning scheme is a concrete instance of Chapter~\ref{ch:ch13-verification-inheritance}'s claim-type policy layer, applied to a model artifact rather than a scholarly claim: format, contract, and content versions are, respectively, an admissibility check, a compatibility-regime check, and a provenance identifier, exactly the three roles that chapter kept deliberately separate.

\section{Channels as Streams of Executable Models}

Typed channels normally transport values between concurrent processes. If a channel value is a model capsule, the receiving process can validate it, allocate a fresh arena, bind inputs, and execute the declared graph. The channel then represents a stream of parameterized computations. This is best understood as mobile model description, not mobile authority: receipt of a capsule must not grant filesystem, network, process, or device access.
\[
\mathrm{send}(C) \to \mathrm{receive}(C) \to \mathrm{verify}(C) \to \mathrm{allocate}(A) \to \mathrm{execute}(C,A,x) \to y.
\]
Fresh arenas provide useful isolation between invocations, eliminating accidental reuse of mutable intermediate state unless it is explicitly declared and transferred. They do not by themselves prevent denial of service, side channels, malformed-shape attacks, or vulnerabilities in the verifier and kernels.

Let a receiver state be $\sigma = (A,M,V)$, where $A$ is the available arena, $M$ the admitted capsule set, and $V$ the runtime version. The transition $\mathrm{Exec}(C,x)$ is admissible only if the capsule signature and hashes verify, its declared operators are supported by $V$, its maximum memory and instruction budget fit $A$, and its input contract accepts $x$. Execution occurs in a new arena $A'$ and returns either a typed result $y$ with an execution record $e$, or a typed refusal $r$:
\[
(\sigma,C,x) \to (\sigma,y,e) \qquad \text{or} \qquad (\sigma,C,x) \to (\sigma,r).
\]
Refusal is part of the protocol, not an exceptional accident, in precisely the sense Chapter~\ref{ch:ch11-constitutional-space-of-refusal} argued refusal must be a first-class, non-punitive outcome rather than an absence of production. The execution record should include capsule hash, runtime hash, device class, input hash or privacy-preserving reference, start and end times, resource counters, and result status; sensitive inputs and outputs need not be placed in the record.

\section{Continuous Data Loops and Feedback Hazards}

Repeated training on self-generated text can narrow linguistic and conceptual diversity, amplify systematic errors, leak evaluation material, and reward artifacts of the validator. A robust loop therefore keeps an immutable external evaluation set, retains real or independently sourced data, measures duplication and source concentration, records ancestry between examples, and limits the number of generations through which unsupported claims can propagate:
\[
D_{t+1} = R_t \cup H_t \cup \{ z \in G_t : V_t(z) = \mathrm{accept} \},
\]
where $R_t$ denotes independently sourced material, $H_t$ human-reviewed examples, $G_t$ generated candidates, and $V_t$ the versioned validator. The decomposition makes it possible to report performance by data origin rather than obscuring the synthetic fraction in a single aggregate.

The scientifically defensible version of an open-versus-private service uses the same training protocol, acceptance tests, and reporting schema in both modes. The variable being purchased is isolation and operational control, not a higher standard of evidence. Public runs may publish data, configuration, metrics, and artifacts; private runs keep customer data and outputs inside a customer-controlled enclave while emitting only the attestations and reports the customer authorizes. Bring-your-own-cloud deployment reduces custody by allowing the customer to own accounts, encryption keys, logs, storage, and teardown, but it does not eliminate trust: the customer must still assess the supplied code, binary provenance, infrastructure declarations, dependencies, and operator access paths.

\section{Security and the Trusted Computing Base}

The relevant adversaries include a malicious capsule producer, a compromised build dependency, an operator with excessive cloud privileges, a customer dataset containing secrets or prompt injections, and an external party observing timing or resource use. The trusted computing base should therefore be as small as practical: parser, verifier, allocator, operator kernels, signature implementation, and the mechanism enforcing resource and I/O policy. Malformed capsules are controlled by canonical parsing, shape validation, and fuzzing, with residual risk in verifier or kernel vulnerability; arbitrary code execution is controlled by restricted bytecode with no native extensions, with residual risk in bugs among the permitted operators; resource exhaustion is controlled by static bounds and runtime quotas, with residual risk in adversarial worst-case scheduling; training-data leakage is controlled by access controls, minimization, and memorization tests, with residual risk that inference-time extraction remains possible; supply-chain substitution is controlled by a pinned closure, hashes, signatures, and reproducible builds, with residual risk in a compromised trusted source or key; and cross-job contamination is controlled by fresh arenas and tenant isolation, with residual risk in hardware and timing side channels. Privacy is a property of the complete workflow, not a synonym for running in a private account; it requires an explicit data-retention policy, logging policy, key ownership model, access review, egress controls, incident procedure, and deletion verification, and neither differential privacy nor confidential-computing hardware should be implied merely by the word ``enclave.''

Shrinking the trusted computing base is only useful to the extent that the cost of exercising it stays bounded, so a deployment claim should report the verification overhead alongside the security boundary. Let $c_v$ denote the cost of the parse-and-verify step for a capsule and $c_x$ the cost of subsequently executing it once. The relative verification overhead $c_v/(c_v+c_x)$ is large and largely irrelevant for a capsule executed many times after a single verification, but it dominates the cost of a workload streaming many distinct small capsules through a receiver, verifying each once before a single execution; the two regimes should not be reported under one aggregate latency figure. A related and separable cost is re-verification triggered by contract-version changes: a runtime upgrade that narrows its accepted contract-version range can force re-verification or outright refusal of a large previously admitted capsule population, an operational cost of the versioning discipline that should be planned for rather than discovered at upgrade time.

\section{Experimental Program and Falsifiable Claims}

The architecture should be tested through preregistered comparisons rather than demonstrations alone. The primary baseline is a conventional floating-point training and post-training export pipeline; the deployment-native condition uses the target quantizer and runtime-equivalent forward semantics during training, with both receiving matched data, parameter counts, optimization budgets, and evaluation procedures. Quality and efficiency must be reported jointly: a speedup measured at lower accuracy, shorter context, a different sampler, or altered batch size is not a controlled comparison, and energy claims require wall-plug or device-level measurement rather than arithmetic-operation counts alone.

A useful experimental sequence first verifies byte-level capsule stability and operator conformance on small networks, then measures deployment disagreement for individual layers and complete models, then compares ternary-native training with post-training quantization, and only after those controls pass adds capsule streaming, multi-tenant execution, and the continuous data loop, an ordering that prevents a complex integration from concealing a basic numerical mismatch. Because $\delta_{\mathrm{out}}$ and $\delta_{\mathrm{num}}$ are estimated from a finite evaluation set, a reported reduction in either quantity should be accompanied by an interval estimate and a preregistered minimum detectable effect, not a point estimate alone, and a family of ablations tested against a shared evaluation set should disclose the resulting multiple-comparisons burden rather than reporting each result as though it were independently significant.

The proposal should be rejected or revised if deployment-native training does not reduce output disagreement relative to a carefully tuned export baseline; if claimed efficiency disappears at matched task quality; if the silver-label loop improves its internal validator while degrading external evaluation; if capsule verification cannot bound resource use; or if private operation depends on undeclared operator access. Negative results would still be informative because they would locate the cost of representation alignment and identify which integration layers provide no measurable benefit, precisely the falsifiability discipline Chapter~\ref{ch:ch02-accommodation-before-prediction} argued a research program needs and a flat-likelihood hypothesis lacks.

\section{Conclusion}

The strongest idea in this architecture is not that ternary cells are miniature intelligent agents or that channels become fundamentally new objects. It is that a deployed neural computation can be packaged as a constrained, typed, verifiable program whose discrete representation is present during optimization, evaluation, transport, and execution. This changes the unit of reproducibility from a weight file to a complete computational contract, and it is the material instance, at the level of an actual deployed system, of the same discipline this Part has developed at every other level: CLIO's separation of documentary history from operative state, the verify seam's separation of consequence from rationale, and Claimshift's separation of a claim's identity from its correspondence under transformation all recur here as the separation of a model's packed weights from the complete deployment specification $F$ that determines what those weights actually compute.
